\chapter{Student inference with unknown variance}
\label{sec:p15-student-inference}

Neyman's Example III constructs confidence limits for a normal
mean when the variance is unknown
\cite[pp.~376--378, equations~(121)--(130)]{Neyman1937Confidence}.
The one-sided comparison comes from the similar-region construction
of Neyman and Pearson
\cite[pp.~313--321, Examples~7--8]{NeymanPearson1933Tests}.
The two-sided local comparison is the normal specialization of
Neyman's composite-hypothesis type-B criterion
\cite[pp.~251--265]{Neyman1935CompositeHypotheses}.
Here \(\alpha\) is the confidence coefficient and
\(\varepsilon=1-\alpha\) is the rejection probability.
The development proves both that original local comparison and
a separate global unbiased comparison.

\section{The model and the Student law}

\begin{hbdefinition}{def:p15-student-model}{Unknown mean and unknown positive variance}
Fix an integer \(n\geq2\), put \(\nu=n-1\), and let
\(X_1,\ldots,X_n\) be independent with law \(N(\mu,\sigma^2)\),
where \(\mu\in\mathbb R\) and \(\sigma>0\) are both unknown.
Write \(P_{\mu,\sigma}\) and \(E_{\mu,\sigma}\) for probability
and expectation on the Borel sample space \(\mathbb R^n\). Define
\[
 \bar X=\frac1n\sum_{i=1}^nX_i,\qquad
 V=\sum_{i=1}^n(X_i-\bar X)^2,\qquad
 S^2=\frac V\nu,\qquad e=\frac S{\sqrt n}
                       =\sqrt{\frac{V}{n\nu}}.
\]
Thus \(S\) is the sample standard deviation and \(e\) is its
associated standard error. Neyman's \(s\) in his 1937
equation~(123) is \(e\). For a candidate mean \(m\), the
Student statistic, on \(V>0\), is
\[
                         T_m=\frac{\bar X-m}{e}.
\]
All tests below are defined on \(V=0\) by their stated unscaled
inequalities, and all confidence intervals include their endpoints.
This convention defines every sample outcome without dividing by zero.
The normal and gamma parameter conventions are those of
\cref{def:p15-normal-law,def:p15-gamma-density}; in particular
the second gamma parameter is a rate.
\index{Student distribution}\index{confidence interval!Student}
\end{hbdefinition}

\begin{hbtheorem}{thm:p15-student-orthogonal}{proved=annex:p15.student.orthogonal}{Orthogonal coordinates and independent residual variance}
There is an orthogonal matrix \(O\), with first row
\(n^{-1/2}(1,\ldots,1)\), for which
\[
 Y=O(X-\mu\mathbf1)/\sigma
\]
has independent standard normal coordinates. Specifically, with
\(b=(1,0,\ldots,0)^{\mathsf T}\),
\(d=n^{-1/2}\mathbf1\), and \(v=b-d\), one may use
\(O=I-2vv^{\mathsf T}/(v^{\mathsf T}v)\).
Then
\[
 Y_1=\frac{\sqrt n(\bar X-\mu)}{\sigma},\qquad
 W=\sum_{j=2}^nY_j^2=\frac V{\sigma^2}.
\]
The variables \(Y_1\) and \(W\) are independent,
\(W\) has law \(\mathrm{Gamma}(\nu/2,1/2)\), and \(V>0\)
almost surely for every \((\mu,\sigma)\).
The notation \(\chi^2_k\), for a positive integer \(k\), means
\(\mathrm{Gamma}(k/2,1/2)\).
\end{hbtheorem}

\begin{hblemma}{thm:p15-student-beta-direction}{proved=annex:p15.student.beta.direction}{Gamma ratios and the exact spherical coordinate}
If \(A,B\) are independent gamma variables with shapes \(a,b>0\)
and common rate \(\lambda>0\), then
\[
 A+B\ \hbox{and}\ \frac A{A+B}
\]
are independent, with respective laws
\(\mathrm{Gamma}(a+b,\lambda)\) and \(\mathrm{Beta}(a,b)\).
Let \(\omega_n\) denote the normalized sphere measure
\(\sigma_n\) defined in \cref{def:p15-wce-spherical-measure}.
For \(U\) with law \(\omega_n\), the projection
\(Z=d^{\mathsf T}U\), \(d=n^{-1/2}\mathbf1\), has density
\[
 h_n(z)=
 \begin{cases}
 \displaystyle
 \frac{\Gamma(n/2)}{\sqrt\pi\,\Gamma((n-1)/2)}
 (1-z^2)^{(n-3)/2},&-1<z<1,\\
 0,&\text{otherwise}.
 \end{cases}
\]
It is symmetric, atomless, and positive on \((-1,1)\).
For independent standard normal \(Y_j\),
\(Y/\|Y\|\) has law \(\omega_n\), independently of
\(\|Y\|^2\sim\chi_n^2\).
\end{hblemma}

\begin{hbtheorem}{thm:p15-student-law}{proved=annex:p15.student.law}{The normalized Student pivot}
For independent \(Z_0\sim N(0,1)\) and \(W\sim\chi_\nu^2\),
the variable \(Z_0/\sqrt{W/\nu}\) has density
\[
 f_\nu(t)=
 \frac{\Gamma((\nu+1)/2)}
      {\sqrt{\nu\pi}\,\Gamma(\nu/2)}
 \left(1+\frac{t^2}{\nu}\right)^{-(\nu+1)/2},
 \qquad t\in\mathbb R.
\]
Its CDF \(F_\nu\) is continuous, strictly increasing and symmetric:
\(F_\nu(-t)=1-F_\nu(t)\). For \(0<p<1\), define
\(t_{\nu,p}\) by \(F_\nu(t_{\nu,p})=p\).
Under every \(P_{\mu,\sigma}\), the pivot \(T_\mu\)
has this same law. Equivalently, for the spherical coordinate above,
\[
        \frac{\sqrt\nu Z}{\sqrt{1-Z^2}}
                \quad\hbox{has CDF }F_\nu.
\]
\end{hbtheorem}

\begin{hbtheorem}{thm:p15-student-coverage}{proved=annex:p15.student.coverage}{Exact one-sided and central confidence intervals}
For every \(0<\alpha<1\), define
\[
\begin{aligned}
 C_-(X)&=[\,\bar X-t_{\nu,\alpha}e,\ \infty),\\
 C_+(X)&=(-\infty,\ \bar X+t_{\nu,\alpha}e\,],\\
 C_0(X)&=[\,\bar X-t_{\nu,(1+\alpha)/2}e,\
                  \bar X+t_{\nu,(1+\alpha)/2}e\,].
\end{aligned}
\]
Each has exact coverage \(\alpha\) for every
\(\mu\in\mathbb R,\sigma>0\).
The inverted rejection tests for a fixed candidate \(m\) are
\[
\begin{aligned}
 \phi_m^+(x)&=\mathbf1_{\{\bar x-m>t_{\nu,\alpha}e(x)\}},\\
 \phi_m^-(x)&=\mathbf1_{\{\bar x-m<-t_{\nu,\alpha}e(x)\}},\\
 \phi_m^0(x)&=\mathbf1_{\{|\bar x-m|>
                             t_{\nu,(1+\alpha)/2}e(x)\}}.
\end{aligned}
\]
These inequalities define the tests even when \(V(x)=0\), where
the intervals become \([\bar x,\infty)\), \((-\infty,\bar x]\),
and \(\{\bar x\}\), respectively.
\end{hbtheorem}

Neyman's one-sided table instruction uses the two-tail table
probability \(2(1-\alpha)\), and his central interval uses
\(1-\alpha\). The signed quantiles above also cover
\(\alpha\leq1/2\), for which the one-sided critical value is
nonpositive. The 1933 paper uses a different scale:
its \(s^2=V/n\) and its \(z=(\bar X-m)/s\), so its
\(z=T_m/\sqrt\nu\). Its integral
\(\int_{z_0}^{\infty}(1+z^2)^{-n/2}\,dz\)
therefore agrees with the one-sided Student quantile after this
change of variable.

\section{Similarity, completeness, and conditional comparisons}

\begin{hbdefinition}{def:p15-student-comparison-classes}{The four testing comparison classes}
Fix the candidate \(m\), and let a randomized test be a Borel
function \(\phi:\mathbb R^n\to[0,1]\). Its power is
\(\beta_\phi(\mu,\sigma)=E_{\mu,\sigma}\phi(X)\).
For \(0<\varepsilon<1\), it has \emph{level}
\(\varepsilon\) when
\(\beta_\phi(m,\sigma)\leq\varepsilon\) for every \(\sigma>0\).
It is \emph{similar of size} \(\varepsilon\) when equality holds
for every \(\sigma>0\).
It is \emph{stationary similar} when it is similar and
\[
           \partial_\mu\beta_\phi(m,\sigma)=0
                    \quad\hbox{for every }\sigma>0.
\]
It is \emph{globally unbiased of level} \(\varepsilon\) for the
two-sided alternative when it has that level and
\[
      \beta_\phi(\mu,\sigma)\geq\varepsilon
           \quad(\mu\ne m,\ \sigma>0).
\]
A uniformly most powerful test in a specified class has power
at least that of every member at every specified alternative.
The abbreviation UMPU means uniformly most powerful among the
globally unbiased level-\(\varepsilon\) tests.

Neyman's composite type-B local criterion requires similarity,
zero first derivative, and the largest possible null second
derivative of power, separately for every nuisance value, among
tests satisfying those first two requirements. In the acceptance
formulation of his 1937 equations~(126)--(129), the second
derivative is instead minimized. This definition makes no
assertion about the realized length of an interval.
\end{hbdefinition}

\begin{hblemma}{thm:p15-student-laplace-completeness}{proved=annex:p15.student.laplace.completeness}{The completeness premise for every nuisance scale}
Let \(a>0\), and let \(g:(0,\infty)\to\mathbb R\) be
measurable with
\[
 \int_0^\infty |g(q)|q^{a-1}e^{-\lambda q}\,dq<\infty
                       \quad\hbox{for every }\lambda>0.
\]
If the same integrals without the absolute value vanish for
every \(\lambda>0\), then \(g=0\) Lebesgue almost everywhere.
Consequently the family
\(\{\mathrm{Gamma}(a,\lambda):\lambda>0\}\) is complete
for every function integrable under every member: zero
expectation under the whole family forces zero almost everywhere.
The proof uses the actual Bernstein approximation of
\cref{thm:p15-feller-bernstein} and establishes the needed signed
Laplace uniqueness explicitly.
\end{hblemma}

\begin{hbtheorem}{thm:p15-student-null-fibers}{proved=annex:p15.student.null.fibers}{Null spheres and the alternative tilt}
For the fixed candidate \(m\), put
\[
 Q_m=\sum_i(X_i-m)^2,\qquad
 U_m=\frac{X-m\mathbf1}{\sqrt{Q_m}},\qquad
 Z_m=d^{\mathsf T}U_m=\frac{\sqrt n(\bar X-m)}{\sqrt{Q_m}}
\]
on \(Q_m>0\); assign an arbitrary unit direction at \(Q_m=0\).
Under \(P_{m,\sigma}\), the variables \(Q_m,U_m\) are
independent with laws
\(\mathrm{Gamma}(n/2,1/(2\sigma^2))\) and \(\omega_n\).
For \(\delta=\mu-m\), write
\[
 a(q)=\frac{\delta\sqrt{nq}}{\sigma^2},\qquad
 M(a)=\int_{S^{n-1}}e^{a d^{\mathsf T}u}\,d\omega_n(u).
\]
Under \(P_{m+\delta,\sigma}\), the conditional direction law
at radius square \(q>0\) has density
\(e^{a(q)d^{\mathsf T}u}/M(a(q))\) relative to \(\omega_n\).
The density of \(Q_m\) is its null gamma density multiplied by
\[
                  e^{-n\delta^2/(2\sigma^2)}M(a(q)).
\]
All these densities are strictly positive on their stated domains.
On \(V>0\),
\[
        T_m=\frac{\sqrt\nu Z_m}{\sqrt{1-Z_m^2}}.
\]
\end{hbtheorem}

\begin{hbtheorem}{thm:p15-student-score-fibers}{proved=annex:p15.student.score.fibers}{Exactly what similarity and unbiasedness force}
Every bounded test has a twice continuously differentiable
power function in \(\mu\), for each fixed \(\sigma>0\).
With \(H_m=\sum_i(X_i-m)\),
\[
\begin{aligned}
 \partial_\mu\beta_\phi(m,\sigma)
     &=E_{m,\sigma}\!\left[\frac{H_m}{\sigma^2}\phi\right],\\
 \partial_\mu^2\beta_\phi(m,\sigma)
     &=E_{m,\sigma}\!\left[
          \left(\frac{H_m^2}{\sigma^4}-\frac n{\sigma^2}\right)\phi
                       \right].
\end{aligned}
\]
Define the Borel angular integrals
\[
\begin{aligned}
 A_\phi(q)&=\int\phi(m\mathbf1+\sqrt q\,u)\,d\omega_n(u),\\
 B_\phi(q)&=\int (d^{\mathsf T}u)
                   \phi(m\mathbf1+\sqrt q\,u)\,d\omega_n(u).
\end{aligned}
\]
Then similarity is equivalent to
\(A_\phi(q)=\varepsilon\) for Lebesgue almost every \(q>0\).
Stationary similarity is equivalent to this equation together
with \(B_\phi(q)=0\) for Lebesgue almost every \(q>0\).
Every globally unbiased level-\(\varepsilon\) test is stationary
similar. The common Lebesgue null sets in these assertions are
null under every radius law in
\cref{thm:p15-student-null-fibers}.
The angular constraints need not hold at every individual radius.
\end{hbtheorem}

\begin{hbtheorem}{thm:p15-student-one-sided-optimal}{proved=annex:p15.student.one.sided.optimal}{The one-sided Student tests are best among similar tests}
The test \(\phi_m^+\) of
\cref{thm:p15-student-coverage} is similar of size
\(\varepsilon=1-\alpha\) and has power at least that of every
similar test at every \(\mu>m,\sigma>0\).
The reflected test \(\phi_m^-\) has the corresponding property
at every \(\mu<m,\sigma>0\).
On each null sphere the upper rejection set is the cap
\[
 Z_m>c_+,\qquad
 c_+=\frac{t_{\nu,\alpha}}
                {\sqrt{\nu+t_{\nu,\alpha}^2}},
\]
and the lower set is \(Z_m<-c_+\).
These are cones about the diagonal direction in sample space;
they impose no requirement that every coordinate have the
same sign. Among jointly measurable confidence procedures
with exact coverage \(\alpha\) for every \((\mu,\sigma)\),
\(C_-\) minimizes inclusion probability of a fixed false
candidate \(m<\mu\), and \(C_+\) does so for \(m>\mu\).
\end{hbtheorem}

\section{The central interval: global and local optimality}

\begin{hblemma}{thm:p15-student-exponential-chord}{proved=annex:p15.student.exponential.chord}{The complete conditional two-sided comparison}
On \(S^{n-1}\) with measure \(\omega_n\), put
\(Z(u)=d^{\mathsf T}u\), which has the symmetric density \(h_n\).
Let
\(c\in(0,1)\) satisfy \(P(|Z|>c)=\varepsilon\).
On the sphere write \(r(u)=\mathbf1_{\{|d^{\mathsf T}u|>c\}}\).
If a measurable \(g:S^{n-1}\to[0,1]\) satisfies
\[
                  \int g\,d\omega_n=\varepsilon,\qquad
                  \int (d^{\mathsf T}u)g\,d\omega_n=0,
\]
then for every real \(a\),
\[
                  \int r e^{aZ}\,d\omega_n
                         \geq\int g e^{aZ}\,d\omega_n.
\]
Moreover, for \(a\ne0\),
\[
                  \frac{\int r e^{aZ}\,d\omega_n}
                       {\int e^{aZ}\,d\omega_n}>\varepsilon.
\]
Under the sole mass constraint \(\int g=\varepsilon\),
\(\int rZ^2\,d\omega_n\geq\int gZ^2\,d\omega_n\).
\end{hblemma}

\begin{hbtheorem}{thm:p15-student-two-sided-optimal}{proved=annex:p15.student.two.sided.optimal}{Global unbiased optimality and Neyman's exact local criterion}
The central Student test \(\phi_m^0\) is stationary similar
of size \(\varepsilon=1-\alpha\), and is strictly unbiased:
\[
             E_{\mu,\sigma}\phi_m^0>\varepsilon
                  \quad(\mu\ne m,\ \sigma>0).
\]
At every such alternative it dominates every stationary
similar test, and in particular it is UMPU.
Its null second derivative of power is at least that of
every similar size-\(\varepsilon\) test, separately for each
\(\sigma>0\). It therefore satisfies the composite type-B
criterion. Its acceptance probability has first derivative
zero and the smallest second derivative in the class of
similar, stationary acceptance procedures.
The central spherical threshold is
\[
 c=\frac{t_{\nu,(1+\alpha)/2}}
              {\sqrt{\nu+t_{\nu,(1+\alpha)/2}^2}}\in(0,1).
\]
\end{hbtheorem}

\begin{hbtheorem}{thm:p15-student-unbiased-confidence}{proved=annex:p15.student.unbiased.confidence}{What the short-unbiased confidence conclusion means}
Allow a confidence procedure to have a jointly Borel inclusion
function \(a(x,m)\in[0,1]\); an ordinary confidence set has
\(a(x,m)=\mathbf1_{\{m\in C(x)\}}\). Call it unbiased at
coefficient \(\alpha\) when
\[
 E_{\mu,\sigma}a(X,\mu)\geq\alpha,\qquad
 E_{\mu,\sigma}a(X,m)\leq\alpha
                       \quad(m\ne\mu,\ \sigma>0).
\]
Then its true coverage is necessarily exactly \(\alpha\).
For every false \(m\), the central Student interval satisfies
\[
 P_{\mu,\sigma}\{m\in C_0(X)\}
                     \leq E_{\mu,\sigma}a(X,m).
\]
It is itself unbiased and has exact coverage.
For each fixed candidate, it also minimizes the null second
derivative of its inclusion probability among procedures
with nuisance-uniform exact coverage and zero first derivative.
These are false-coverage and local-curvature comparisons;
no minimum of the length of each realized interval is asserted.
The coverage probability concerns repeated sampling from
each fixed \((\mu,\sigma)\) and supplies no posterior law
for \(\mu\) or \(\sigma\).
\end{hbtheorem}

\begin{hblemma}{thm:p15-student-source-score-bridge}{proved=annex:p15.student.source.score.bridge}{The normal specialization of Neyman's type-B construction}
Use source parameters \(\theta_1=\mu\), \(\theta_2=v=\sigma^2\).
At the candidate \(m\), the log-density scores and their
derivatives are
\[
\begin{aligned}
 \varphi_1&=H_m/v,&
 \varphi_2&=-n/(2v)+Q_m/(2v^2),\\
 \varphi_{11}&=-n/v,&
 \varphi_{12}&=-\varphi_1/v,&
 \varphi_{22}&=-n/(2v^2)-2\varphi_2/v.
\end{aligned}
\]
They satisfy the affine derivative relations of the 1935
construction. Fixing its nuisance score is exactly fixing
\(Q_m\). Its conditional size and first-score constraints
are the two angular equations of
\cref{thm:p15-student-score-fibers}; its remaining quadratic
objective is a positive multiple of the \(Z_m^2\) objective
in \cref{thm:p15-student-exponential-chord}.
The completeness and differentiation premises needed for this
specialization hold by the preceding proofs.
\end{hblemma}

\section{Small samples and boundaries}

\begin{hbtheorem}{thm:p15-student-worked-cases}{proved=annex:p15.student.worked.cases}{Exact calculations for two and three observations}
For \(\nu=1\),
\[
 f_1(t)=\frac1{\pi(1+t^2)},\qquad
 F_1(t)=\frac12+\frac{\arctan t}{\pi},\qquad
 t_{1,p}=\tan\!\left(\pi(p-\tfrac12)\right).
\]
For observations \((1,3)\), \(\bar x=2,V=2,e=1\).
The one-sided intervals of confidence \(3/4\) are
\([1,\infty)\) and \((-\infty,3]\).
The central interval of confidence \(1/2\) is \([1,3]\);
that of confidence \(3/4\) is
\([1-\sqrt2,\,3+\sqrt2]\).

For \(\nu=2\),
\[
 F_2(t)=\frac12+\frac{t}{2\sqrt{2+t^2}},\qquad
 t_{2,p}=\frac{\sqrt2(2p-1)}{\sqrt{1-(2p-1)^2}}.
\]
For \((0,1,2)\), \(\bar x=1,V=2,e=1/\sqrt3\).
The central interval of confidence \(9/10\) is
\[
                   [\,1-3\sqrt{6/19},\ 1+3\sqrt{6/19}\,].
\]
The pivot with \(\nu=1\) is symmetric but has no finite mean.
For \(n=1\) the stated residual scale is unavailable.
The independence assertion depends on the sampling law:
with two independent fair
\(\{-1,1\}\)-valued observations these two statistics are
dependent. Even under the normal model, replacing a Student
critical value by a normal critical value changes coverage;
for \(n=2\), an interval \(\bar X\pm c e\), \(c>0\), has
coverage \(2\arctan(c)/\pi\).
\end{hbtheorem}
